\section{The construction side}\label{sec:construction}

This section says what the two theorems exclude, what they leave, and what would have to be added to make them
statements about machines. Readings of the one-hop model as statements about tapes are labelled as readings.

\subsection{Mechanisms that fail}\label{sec:fail}

\Cref{tab:fail} lists upper bounds on the code dimension for decoders with structure on one index, in the linear
one-hop model, whatever is done on the other index. Each row is a theorem of \cref{sec:order,sec:symmetry}.

\begin{table}[ht]
\centering\small
\begin{tabular}{@{}>{\raggedright\arraybackslash}p{0.47\textwidth}>{\raggedright\arraybackslash}p{0.23\textwidth}l@{}}
\toprule
Structure on one index & Upper bound & By \\
\midrule
At most two Kronecker terms & $d(d+1)/2$ & \cref{cor:two-terms} \\
Commuting maps: any convolutions of any number of rows read in step & $d(d+1)/2$ & \cref{cor:commutative} \\
One $p$-group in characteristic~$p$: one dyadic hierarchy over $\F_2$ & $d(d+1)/2$ & \cref{cor:commutative} \\
All affine index maps, $d=p^m$, characteristic $p$ & $d(d+1)/2$ & \cref{cor:affine} \\
Expansion factor $c$ along some flag & $(1-\frac1{2c})d^2+O(d)$ & \cref{cor:expansion} \\
Uniform decoders with $M$ interval pieces, cycle or line & $(d^2+Md)/2$ & \cref{thm:TP,thm:line} \\
Copy rules with $M$ decoder classes of any shape & $(d^2+Md)/2$ & \cref{prop:copy} \\
$s$ modulation frequencies & $\frac s{s+1}(d^2+d)$ & \cref{cor:weyl} \\
One address digit of size $q$ treated uniformly & $\frac12(d^2+d^2/q)$ & \cref{cor:digit} \\
\bottomrule
\end{tabular}
\caption{Upper bounds on the code dimension in the linear one-hop model. The other index is arbitrary.}
\label{tab:fail}
\end{table}

Two design rules follow, for the one-hop model.
\begin{itemize}[leftmargin=1.5em]
\item \emph{Synchronous access is the pairing.} Reading any number of source rows in step with the output, at any
offsets and through any convolution, gives the symmetric matrices and no more.
\item \emph{Uniform passes are the pairing.} A stretch on which translated outputs share their decoder gains nothing,
whatever its window. The bound rises only with changes of regime, by $d/2$ symbols for each.
\end{itemize}

\subsection{A reading for tapes}\label{sec:reading}

The following is a reading, not a theorem about machines. A pass in which the position read and the position written
advance at the same speed, by a linear finite-state rule, computes a convolution with an unbounded window. On a line
this is a two-sided rule in the sense of \cref{thm:line}, with one class per stretch of unchanged regime, under two
provisos: the regime changes at the same positions for every output row (change points that depend on the row need
their common refinement), and the rule of each regime is applied to the whole stored word, with no state carried
across a change. A state of $q$ symbols carried across a change alters $K$ by a rank of at most $q$ per output row, so
by \cref{lem:pigeonhole} it costs a further term of order $qd$ per change. Under the provisos, \cref{thm:line} bounds
the gain over the pairing value by $d/2$ symbols per regime change, so rate $\tfrac12+\delta$ needs at least $2\delta d$
regime changes, and for equal pieces the bound is attained by the block-symmetric matrices, that is, by the coarse step
of the folklore recursion. Other speed ratios are decimations. On a cycle a single ratio is one twist class of
\cref{thm:twist}; mixtures of ratios fall under \cref{cor:expansion} or \cref{cor:affine}.

This reading matches what the series proved about motion (\cref{sec:motion}): locally affine motion is excluded on the
geometric side by the second and fourth papers, and here uniform decoding is excluded on the coding side, in a model
that has no geometry at all. The seventh paper has a statement of the same kind for storage codes on relations: Cayley
relations $\mathrm{Cay}(G,\{s,s^{-1},t\})$ without transitive triangles give no gain over the pairing rate, for every
finite group~\cite[Corollary~5.12]{SeventhPaper2026}. \Cref{thm:TP} allows any degree and any window, but only one hop and only linear
decoders.

\subsection{What a linear-time machine would need}\label{sec:would-need}

Combining \cref{sec:needs} with \cref{tab:fail}: in addition to positive characteristic, curved motion, expanding cone
relations and growing paired distance, a linear-time transposer, read in the one-hop model, needs on \emph{both}
indices a family of maps that
\begin{itemize}[leftmargin=1.5em]
\item has no invariant flag and no pair of flags with bounded expansion (\cref{thm:FL});
\item commutes with no non-derogatory translation up to $o(d)$ boundary columns, and intertwines none with boundedly
many pairwise commuting others (\cref{thm:TP,thm:twist});
\item treats no address digit uniformly (\cref{cor:digit}).
\end{itemize}
The digit masks at every scale, which the folklore recursion uses, have these properties: their number grows with
$d$, and no shift respects them. A bounded number of classes is not enough for rate near one, whatever their shape,
because $M$ classes leave a flag with $g(i)\le Mi$ and hence rate at most $1-1/(2M)+O(1/d)$. What the theorems leave
open is how large a code a given number of classes can carry when the classes have no almost-period. It can exceed the
value $(d^2+Md)/2$ of every copy rule and of every interval rule, at $d=6$ though not at $d=5$
(\cref{prop:aperiodic,rem:small}); whether a gain survives in rate is \cref{prob:aperiodic}.

\subsection{From the one-hop model to machines}\label{sec:to-machines}

Three steps separate \cref{thm:TP} from a statement about oblivious machines.
\begin{enumerate}[leftmargin=2em]
\item \emph{A hypothesis on cone systems.} The coding bounds of the series use properties of the cone relation, and
\cref{sec:false} shows that the relation-level properties proved so far cannot carry the conclusion. ``Translated receivers share their
decoder on every stretch of periodic head motion'' is a property of the rule. It has to be stated for the cone systems
of real schedules, where the receivers are output bits and the decoders are compositions of the transition function.
\item \emph{Depth reduction.} The genie of \cref{sec:depth} reveals values at chosen steps. The revealed word is advice
that depends on position, so it breaks uniformity wherever it is used. A version of \cref{thm:TP} with advice would have
to charge each revealed bit as a regime change; \cref{thm:TP} suggests the exchange rate, half a row per change, but
proves nothing about it.
\item \emph{Many hops.} \Cref{rem:compose-TP} bounds $t$ uniform hops by a loss that grows linearly in $t$. A machine
that runs for $cL$ steps has unboundedly many hops, and \cref{rem:compose-TP} shows that one cheap operation can break
every digit symmetry at once. So an argument for machines has to count something that an unshuffle does not supply for free.
\end{enumerate}
We do not know how to take any of these steps.
